\section*{Sets and Indices}

\begin{itemize}
    \item Product types \(i \in P = \{A,B\}\).
    \item Processes \(j \in R = \{\mathrm{mfg},\mathrm{asm},\mathrm{insp}\}\), representing manufacturing, assembly, and inspection.
\end{itemize}

\section*{Parameters}

All parameter values are taken from the CSV files.

\begin{itemize}
    \item Processing time per unit of product \(i\) on process \(j\): \(a_{ji}\).
    \[
    \begin{aligned}
    &a_{\mathrm{mfg},A}=20,\quad a_{\mathrm{asm},A}=5,\quad a_{\mathrm{insp},A}=3,\\
    &a_{\mathrm{mfg},B}=0,\quad a_{\mathrm{asm},B}=7,\quad a_{\mathrm{insp},B}=6.
    \end{aligned}
    \]
    \item Weekly capacity of process \(j\): \(C_j\).
    \[
    C_{\mathrm{mfg}}=120,\qquad C_{\mathrm{asm}}=80,\qquad C_{\mathrm{insp}}=40.
    \]
    \item Hourly process cost (also used as idle-time weight) for process \(j\): \(w_j\).
    \[
    w_{\mathrm{mfg}}=12,\qquad w_{\mathrm{asm}}=8,\qquad w_{\mathrm{insp}}=10.
    \]
    \item Selling price per unit of product \(i\): \(s_i\).
    \[
    s_A=650,\qquad s_B=725.
    \]
    \item Minimum weekly profit requirement: \(P^{\min}=3000\).
    \item Minimum weekly production of Type A: \(L_A=5\).
\end{itemize}

The code computes unit cost and unit profit as
\[
c_i=\sum_{j\in R} w_j a_{ji},\qquad \pi_i=s_i-c_i.
\]
Numerically,
\[
c_A=20(12)+5(8)+3(10)=310,\qquad \pi_A=650-310=340,
\]
\[
c_B=0(12)+7(8)+6(10)=116,\qquad \pi_B=725-116=609.
\]

\section*{Decision Variables}

\begin{itemize}
    \item Weekly production quantity of Type A: \(x_A\) (code: \texttt{x\_A} in Phase 1, \texttt{x\_A2} in Phase 2), integer.
    \item Weekly production quantity of Type B: \(x_B\) (code: \texttt{x\_B} in Phase 1, \texttt{x\_B2} in Phase 2), integer.
\end{itemize}

Bounds implemented in code:
\[
x_A \ge L_A,\qquad x_B \ge 0,\qquad x_A,x_B \in \mathbb{Z}.
\]

\section*{Objective}

The implemented logic is lexicographic, solved in two sequential optimization phases.

\paragraph{Phase 1: minimize weighted idle time.}
For any feasible \((x_A,x_B)\), idle time on process \(j\) is
\[
I_j = C_j - \sum_{i\in P} a_{ji}x_i.
\]
The Phase-1 problem is
\[
\min \; Z_{\mathrm{idle}}=\sum_{j\in R} w_j I_j
\]
subject to the constraints listed below. Let the optimal value be \(Z_{\mathrm{idle}}^\star\).

Equivalently, with substituted coefficients,
\[
\min \; 12\bigl(120-20x_A-0x_B\bigr)+8\bigl(80-5x_A-7x_B\bigr)+10\bigl(40-3x_A-6x_B\bigr).
\]

\paragraph{Phase 2: maximize profit subject to minimum idle time.}
Using the Phase-1 optimum, the code solves
\[
\max \; Z_{\mathrm{profit}}=\pi_A x_A+\pi_B x_B
\]
subject to the same constraints as in Phase 1 and additionally
\[
\sum_{j\in R} w_j\left(C_j-\sum_{i\in P} a_{ji}x_i\right)\le Z_{\mathrm{idle}}^\star + 10^{-5}.
\]
Numerically, the Phase-2 objective is
\[
\max \; 340x_A+609x_B.
\]

\section*{Constraints}

The following constraints are enforced in both phases.

\paragraph{Capacity constraints}
\[
20x_A+0x_B \le 120 \qquad \text{(manufacturing)}
\]
\[
5x_A+7x_B \le 80 \qquad \text{(assembly)}
\]
\[
3x_A+6x_B \le 40 \qquad \text{(inspection)}
\]

\paragraph{Minimum profit constraint}
\[
340x_A+609x_B \ge 3000
\]

\paragraph{Minimum Type A production}
\[
x_A \ge 5
\]

\paragraph{Nonnegativity and integrality}
\[
x_B \ge 0,\qquad x_A,x_B \in \mathbb{Z}
\]

\paragraph{Phase-2 idle-time restriction}
Only in Phase 2:
\[
12(120-20x_A)+8(80-5x_A-7x_B)+10(40-3x_A-6x_B)\le Z_{\mathrm{idle}}^\star+10^{-5}.
\]

\section*{Notes on Implementation}

\begin{itemize}
    \item The code does \emph{not} optimize a single weighted combination of profit and idle time. It uses a strict two-stage lexicographic procedure:
    \begin{enumerate}
        \item minimize weighted idle time;
        \item among solutions with minimum weighted idle time (up to a tolerance of \(10^{-5}\)), maximize profit.
    \end{enumerate}
    \item The hourly process costs are used twice in the implementation: first to compute per-unit profit coefficients \(\pi_i\), and second as weights in the idle-time objective.
    \item Idle times are not introduced as explicit decision variables in the code; they are affine expressions of \(x_A\) and \(x_B\).
    \item Overtime is disallowed implicitly by the capacity inequalities; there are no overtime variables or penalties.
    \item The model is solved as an integer linear program in each phase using PuLP with the Gurobi command interface (\texttt{pulp.GUROBI\_CMD}).
    \item The reported objective value printed by the script is the Phase-2 optimal profit \(Z_{\mathrm{profit}}\), not the minimized idle-time value.
\end{itemize}
